\section{Normalization and the training procedure}
\label{sec:conference-setting}

Let $x$ be a centered feature vector with $C = \E[xx^\top] \succ 0$, and suppose
that $y = x^\top\beta^* + \varepsilon$, where $\E[\varepsilon\mid x] = 0$ and
$\E[\varepsilon^2\mid x] = \sigma^2$. We divide each feature by a positive scale
factor, collected in the diagonal matrix $S$. Training on $z = S^{-1}x$ with
parameter $\theta = S\beta$ then gives the following SGD update on the original
coefficient scale:
\begin{equation}
  \beta_{k+1}
  = \beta_k
    - \frac{\eta S^{-2}}{B} \sum_{i=1}^B x_{k,i}(x_{k,i}^\top\beta_k - y_{k,i}). \label{eq:conference-update}
\end{equation}
We use the average squared loss with a factor $1/2$. In the population
comparisons, we fix the centered intercept at zero and initialize $\beta_0 = 0$.
Each update uses $B$ independent fresh observations, so $k$ updates require
$M = kB$ observations. In Section~\ref{sec:conference-experiments}, we also
sample with replacement from a fixed training set of size $n$. There, $M$ counts
processed observations and can exceed $n$.

We can compensate exactly for global rescaling by adjusting the learning rate.
Replacing $S$ by $gS$ and $\eta$ by $g^2\eta$ leaves every update unchanged when
the methods use the same data stream. Changing the relative feature scales,
however, generally changes $\eta S^{-2}$ in a way that a single learning rate
cannot undo. Whether this leads to a risk difference after tuning requires a
separate comparison.

Write $e_k = \E[\beta_k-\beta^*]$ and
$V_k = \E[(\beta_k-\beta^*)(\beta_k-\beta^*)^\top]$. Fresh sampling gives
\begin{equation}
  V_k = e_ke_k^\top + A_k + \sigma^2N_k, \label{eq:conference-decomposition}
\end{equation}
where $A_k \succeq 0$ is the covariance of the noiseless process under random
feature sampling and $\sigma^2N_k \succeq 0$ is the covariance contributed by
observation noise. This decomposition separates incomplete learning from the two
sources of variation in training. We derive the exact recurrence in
Appendix~\ref{sec:comparison-proofs}.

For a subgroup $G$, we define the aggregate and subgroup excess prediction risks
using half squared loss:
\begin{equation}
  R_k = \tfrac12\operatorname{tr}(CV_k),\qquad R_{G,k}
      = \tfrac12\operatorname{tr}(C_GV_k), \label{eq:conference-risks}
\end{equation}
Here $C_G = \E[xx^\top\mid G]$ is the conditional second moment, which includes
the contribution of the conditional feature mean. Both risks exclude irreducible
test noise. We evaluate both objectives for the same trained predictor, without
changing the distribution of its training data.

\paragraph{Binary features.}

Let $I$ and $J$ be common and rare binary features, respectively, with
\[
  \begin{gathered}
    \Prob(I = 1) = \tfrac12,\qquad \Prob(J = 1) = q, \\
    \Prob(I = 1,J = 1) = \tfrac{q(1+\ell)}2,
  \end{gathered}
\]
where $0 < q < 1/2$ and $-1 \leq \ell \leq 1$. These ranges ensure that all four
cell probabilities are nonnegative. For $x = (I - 1/2,J - q)^\top$, write
$a = 1/4$, $b = q(1 - q)$, and $c = q\ell/2$. Then
\[
  C = \begin{pmatrix}
        a & c \\
        c & b
      \end{pmatrix},\qquad C_R
    = \begin{pmatrix}
        a             & \ell(1 - q)/2 \\
        \ell(1 - q)/2 & (1 - q)^2
      \end{pmatrix},
\]
where $R$ denotes the subgroup $J = 1$. A rare-coordinate error of one
contributes $(1 - q)^2/2$ to rare risk and $q(1 - q)/2$ to aggregate risk.

We compare unscaled features (U), with $P_U = I_2$, and standard-deviation
scaling (S), with $P_S = \diag(a^{-1},b^{-1})$, where $P = S^{-2}$. In the main
theorem, each rule chooses its own constant learning rate and stopping count.
For the supporting comparisons, we calibrate the rates through the common
coordinate. We use the final iterate in both settings.

\paragraph{Damped scaling.}

The sampled-data experiments also include damped scaling, with
\begin{equation}
  s_j = \sqrt{v_j + 0.01},\qquad P_{jj}
      = \frac{1}{v_j + 0.01}, \label{eq:conference-damped}
\end{equation}
where $v_j$ is either the known population variance or an estimate from the
training data. We use the heuristic offset $0.01$ from the controlled pilot,
fixed in variance units before evaluation. We neither derive this value from
theory nor select it by validation. At the same learning rate, damping
multiplies the SD update on the original coefficient scale by
$v_j/(v_j + 0.01)$. This is approximately one half at binary prevalence $0.01$
and $0.96$ for a balanced binary feature, illustrating how the offset reduces
updates more strongly for rare features. The comparison does not establish an
optimal damping level.
